\section{The genus-two Jacobian}

The genus-\(2\) factor \(\Jac(X_A)\) carries the richest arithmetic information in the decomposition. In this section we record Frobenius data at small good primes, establish the automorphic decomposition of \(L(H^1(X),s)\), obtain a Picard-number lower bound for \(\Jac(X)\), and analyze the local reduction and Galois representations of \(\Jac(X_A)\).

\subsection{Frobenius data at small good primes}

\begin{proposition}\label{thm:simple}
For the genus-two curve
\[
X_A:\quad w^2=-y(y-1)(256y^3+411y^2+165y+32),
\]
the Frobenius characteristic polynomials at the good primes \(p=5,13,17\) are
\begin{align}
\chi_5(T)&=T^4-T^3+4T^2-5T+25=(T^2+2T+5)(T^2-3T+5),\label{eq:chi5-factor}\\
\chi_{13}(T)&=T^4-2T^2+169,\label{eq:chi13}\\
\chi_{17}(T)&=T^4-3T^3+28T^2-51T+289.\label{eq:chi17}
\end{align}
In particular, \(\chi_{13}(T)\) is irreducible over~\(\QQ\). Moreover,
\[
\chi_{13}(T)=(T^2-2\sqrt{7}\,T+13)(T^2+2\sqrt{7}\,T+13)
\]
over \(\QQ(\sqrt{7})\), while
\[
\chi_{17}(T)=\left(T^2-\frac{3+\sqrt{33}}2\,T+17\right)\left(T^2-\frac{3-\sqrt{33}}2\,T+17\right)
\]
over \(\QQ(\sqrt{33})\).

More generally, for any good prime \(p\) write
\[
\chi_p(T)=T^4-s_1(p)T^3+s_2(p)T^2-s_1(p)pT+p^2.
\]
Then \(\chi_p(T)\) is reducible over~\(\QQ\) if and only if either
\[
\Delta_p:=s_1(p)^2-4\bigl(s_2(p)-2p\bigr)
\]
is a square in~\(\QQ\), or \(s_1(p)=0\) and \(-s_2(p)-2p\) is a square in~\(\QQ\). For
\[
p=13,17,23,29,43,47,53
\]
one finds
\[
\Delta_p=112,\ 33,\ 201,\ 73,\ 304,\ 228,\ 112,
\]
and the second case does not occur, so \(\chi_p(T)\) is irreducible over~\(\QQ\) at each of these primes.
\end{proposition}

\begin{proof}
Point counting on \(X_A\) gives
\[
\#X_A(\F_5)=5,\qquad \#X_A(\F_{25})=33,
\]
\[
\#X_A(\F_{13})=14,\qquad \#X_A(\F_{169})=166,
\]
\[
\#X_A(\F_{17})=15,\qquad \#X_A(\F_{289})=337.
\]
Hence
\[
(s_1,s_2)=(1,4)\ \text{at }p=5,\qquad
(s_1,s_2)=(0,-2)\ \text{at }p=13,\qquad
(s_1,s_2)=(3,28)\ \text{at }p=17,
\]
which yields \eqref{eq:chi5-factor}, \eqref{eq:chi13}, and \eqref{eq:chi17}.

For a good prime \(p\), any factorization of
\[
\chi_p(T)=T^4-s_1T^3+s_2T^2-s_1pT+p^2
\]
over~\(\QQ\) cannot have a linear factor: a rational root \(\alpha\) would satisfy \(\alpha^2=p\). Thus a factorization consists of two monic quadratics. Since every Frobenius root has absolute value \(\sqrt p\), each quadratic has constant term \(+p\) or \(-p\), and the two signs agree. In the positive case,
\[
\chi_p(T)=(T^2-uT+p)(T^2-vT+p),
\]
and comparison of coefficients gives
\[
u+v=s_1,\qquad uv=s_2-2p.
\]
This factorization exists over~\(\QQ\) if and only if the quadratic equation
\[
X^2-s_1X+(s_2-2p)=0
\]
has a rational root, equivalently if and only if
\[
\Delta_p=s_1^2-4(s_2-2p)
\]
is a square in~\(\QQ\). In the negative case, comparison of the \(T^3\)- and \(T\)-coefficients forces \(s_1=0\) and opposite linear terms, so
\[
\chi_p(T)=(T^2-uT-p)(T^2+uT-p),\qquad u^2=-s_2-2p.
\]
This gives the stated two-case criterion. The second case requires \(s_1=0\) and \(s_2\leq-2p\); it does not occur at any prime considered here, since the only such prime with \(s_1=0\) is \(p=13\), where \(-s_2-2p=-24<0\). At \(p=5\), one has \(\Delta_5=25\), which yields the displayed factorization. At \(p=13\) and \(p=17\), the values are \(\Delta_{13}=112\) and \(\Delta_{17}=33\), so both quartics are irreducible over~\(\QQ\). The same computation from the point-count data gives
\[
(s_1,s_2)=(3,-2),\ (-1,40),\ (-4,14),\ (6,46),\ (-8,94)
\]
at \(p=23,29,43,47,53\), respectively, hence the additional nonsquare values \(201,73,304,228,112\).

The displayed factorizations over \(\QQ(\sqrt{7})\) and \(\QQ(\sqrt{33})\) are obtained by solving
\[
u+v=s_1,\qquad uv=s_2-2p,
\]
for the quadratic traces \(u,v\) of Frobenius on the two degree-\(2\) factors.
\end{proof}

\begin{corollary}\label{cor:q-simple}
The Jacobian \(\Jac(X_A)\) is simple over~\(\QQ\).
\end{corollary}

\begin{proof}
If not, Poincar\'e reducibility over~\(\QQ\) would give \(\Jac(X_A)\sim_{\QQ}E_1\times E_2\) with elliptic curves \(E_1,E_2\) over~\(\QQ\). Isogenous abelian varieties have isomorphic \(\ell\)-adic representations, so \(E_1\) and \(E_2\) have good reduction at~\(13\), and \(\chi_{13}=\chi_{13}(E_1)\chi_{13}(E_2)\) would be a product of two quadratics in~\(\ZZ[T]\), contradicting the irreducibility in Proposition~\ref{thm:simple}.
\end{proof}

This conclusion does not imply geometric simplicity and does not determine \(\End^0(\Jac(X_A)_{\overline{\QQ}})\).

\begin{remark}\label{rem:frobenius-witnesses}
Corollary~\ref{cor:q-simple} proves that \(\Jac(X_A)\) is simple over~\(\QQ\). The further irreducibility checks at good primes up to~\(53\) provide evidence about the geometric endomorphism algebra, but do not determine it or imply geometric simplicity; they also reinforce the contrast with the explicitly split Prym surface \(P\sim_{\QQ}E_{\mathrm{res}}^2\).
\end{remark}

\subsection{Automorphic structure}

\begin{theorem}\label{thm:automorphic}
The \(L\)-function of the first cohomology of the regular \(S_4\)-closure \(X\) admits the factorization
\begin{equation}\label{eq:automorphic-factor}
L(H^1(X),s)=L(\Jac(X_A),s)\,L(f_{1147},s)^2\,L(f_{37},s)^3\,L(H^1(Q),s)^3,
\end{equation}
where:
\begin{enumerate}
\item[(i)] \(f_{37}\in S_2(\Gamma_0(37))\) is the unique normalized newform of level \(37\), attached to the elliptic curve \(E=37\mathrm{a}1\);
\item[(ii)] \(f_{1147}\in S_2(\Gamma_0(1147))\) is the newform attached to the resolvent elliptic curve \(E_{\mathrm{res}}\), which has conductor \(1147\) and multiplicative reduction at both \(31\) and \(37\);
\item[(iii)] \(L(\Jac(X_A),s)\) is the degree-\(4\) Hasse--Weil \(L\)-function of the genus-\(2\) Jacobian \(\Jac(X_A)\);
\item[(iv)] \(Q=\Prym(Y/E_{\mathrm{res}})\) is an abelian threefold; \(H^1(Q)\) carries the \(V_3^-\)-multiplicity factor.
\end{enumerate}
\end{theorem}

\begin{proof}
Theorem~\ref{thm:full-factor} gives the Hodge-structure decomposition, and the \(L\)-function identity \eqref{eq:automorphic-factor} follows by taking Euler products. Modularity of \(L(E,s)\) and \(L(E_{\mathrm{res}},s)\) is the modularity theorem for elliptic curves over~\(\QQ\) \cite{WilesTaylor1995}. Proposition~\ref{prop:eres-conductor} identifies \(E_{\mathrm{res}}\) as an elliptic curve of conductor \(1147\). The remaining assertions are definitions of the factors appearing in the decomposition.
\end{proof}

\begin{corollary}\label{cor:potential-modularity}
The Hasse--Weil \(L\)-function \(L(\Jac(X_A),s)\) has meromorphic continuation to~\(\mathbb{C}\) and satisfies the expected functional equation.
\end{corollary}

\begin{proof}
Every abelian surface over a totally real field is potentially modular \cite{BCGP2021}; in particular \(L(\Jac(X_A),s)\) has meromorphic continuation and satisfies the expected functional equation.
\end{proof}

\begin{remark}\label{rem:structural-contrast}
The two abelian surface factors of \(\Jac(X)\) are structurally quite different. The sign factor is the Jacobian \(\Jac(X_A)\) of a concrete genus-\(2\) curve, whereas the \(V_2\)-factor \(P\sim E_{\mathrm{res}}^2\) is explicitly split and has endomorphism algebra \(M_2(\QQ)\). Thus the same branch cubic \(c(y)\) produces both a Jacobian-type abelian surface and a non-simple Prym surface.
\end{remark}

\subsection{Explicit endomorphism blocks}

\begin{proposition}\label{thm:non-isogeny}
The elliptic curves \(E\) and \(E_{\mathrm{res}}\) are not \(\QQ\)-isogenous. Moreover,
\[
\End^0(\Jac(X)_{\overline{\QQ}})
\supseteq
\End^0(\Jac(X_A)_{\overline{\QQ}})
\times M_2(\QQ)\times M_3(\QQ)\times \End^0(Q^3_{\overline{\QQ}}).
\]
\end{proposition}

\begin{proof}
For the first assertion, \(\QQ\)-isogenous elliptic curves have the same conductor. Since \(\mathrm{cond}(E)=37\neq 1147=\mathrm{cond}(E_{\mathrm{res}})\), the curves \(E\) and \(E_{\mathrm{res}}\) are not \(\QQ\)-isogenous.

The displayed containment is immediate from the \(\QQ\)-isogeny
\[
\Jac(X)\sim_{\QQ}\Jac(X_A)\times E_{\mathrm{res}}^2\times E^3\times Q^3
\]
of Theorem~\ref{thm:full-factor}. The matrix blocks \(M_2(\QQ)\) and \(M_3(\QQ)\) come from the standard endomorphisms of the powers \(E_{\mathrm{res}}^2\) and \(E^3\).
\end{proof}

\begin{corollary}\label{cor:picard}
The geometric Picard number of \(\Jac(X)\) satisfies
\[
\rho(\Jac(X)_{\overline{\QQ}})\geq 16.
\]
\end{corollary}

\begin{proof}
The isogeny decomposition of Theorem~\ref{thm:full-factor} yields
\[
\rho(\Jac(X)_{\overline{\QQ}})\geq \rho(\Jac(X_A)_{\overline{\QQ}})+\rho(E_{\mathrm{res}}^2{}_{\overline{\QQ}})+\rho(E^3_{\overline{\QQ}})+\rho(Q^3_{\overline{\QQ}}).
\]
Now \(\rho(\Jac(X_A)_{\overline{\QQ}})\geq 1\). For the non-CM elliptic curves \(E_{\mathrm{res}}\) and \(E\), the Rosati-symmetric parts of \(M_2(\QQ)\) and \(M_3(\QQ)\) have dimensions \(3\) and \(6\), so
\[
\rho(E_{\mathrm{res}}^2{}_{\overline{\QQ}})=3,\qquad \rho(E^3_{\overline{\QQ}})=6.
\]
Finally, \(Q^3\) carries the obvious matrix endomorphisms from the three coordinate factors, so its Rosati-symmetric endomorphisms contain the \(6\)-dimensional space of symmetric \(3\times 3\) rational matrices. Hence \(\rho(Q^3_{\overline{\QQ}})\geq 6\), and therefore
\[
\rho(\Jac(X)_{\overline{\QQ}})\geq 1+3+6+6=16.
\]
\end{proof}

\subsection{The Prym threefold}

\begin{proposition}\label{prop:prym-threefold}
The abelian threefold \(Q=\Prym(Y/E_{\mathrm{res}})\) satisfies:
\begin{enumerate}
\item[(i)] \(\Jac(Y)\sim_{\QQ}E_{\mathrm{res}}\times Q\).
\item[(ii)] The \(\ell\)-adic Galois representation on the first cohomology of \(Q\) is
\[
V_\ell(Q)\cong M_{V_3^-}\otimes \QQ_\ell,
\]
where \(M_{V_3^-}\) is the multiplicity space of the anti-standard representation \(V_3^-=V_3\otimes\mathrm{sgn}\) of~\(S_4\).
\item[(iii)] The \(L\)-function of \(Q\) is a degree-\(6\) Euler product. The factorization~\eqref{eq:automorphic-factor} gives
\[
L(H^1(Q),s)^3=\frac{L(H^1(X),s)}{L(\Jac(X_A),s)\,L(E_{\mathrm{res}},s)^2\,L(E,s)^3},
\]
so \(L(H^1(Q),s)\) is determined by the \(L\)-functions of the remaining factors and the genus-\(16\) curve \(X\).
\end{enumerate}
\end{proposition}

\begin{proof}
Part~(i) is the standard Prym decomposition for the degree-\(2\) cover \(Y\to E_{\mathrm{res}}\) \cite{Mumford1974,LangeOrtega2011}.

For~(ii), Theorem~\ref{thm:full-factor} identifies \(H^0(Q,\Omega_Q^1)\) with \(M_{V_3^-}\); the same identification holds at the \(\ell\)-adic level. Part~(iii) follows by dividing \eqref{eq:automorphic-factor} by the three known factors and comparing Euler products at unramified primes.
\end{proof}

\begin{remark}\label{rem:q-limitations}
The threefold \(Q\) is explicit in the sense that it is defined as the Prym variety of the quotient map \(Y=X/C_4\to E_{\mathrm{res}}=X/D_4\), and its cohomological factor is identified through the \(V_3^-\)-multiplicity space. We do not give an affine or projective model for \(Y\), nor do we determine the conductor, endomorphism algebra, or full geometry of \(Q\). The trace table below should therefore be read as a consistency check and arithmetic profile rather than a complete determination of \(Q\).
\end{remark}

\begin{proposition}\label{prop:prym-traces}
For every prime \(p\nmid 2\cdot 3\cdot 31\cdot 37\), the Frobenius trace \(a_1(Q,p)\) is determined by
\[
a_1(Q,p)=\frac{a_1(X,p)-a_1(\Jac(X_A),p)-2\,a_1(E_{\mathrm{res}},p)-3\,a_1(E,p)}{3},
\]
where
\[
\#X(\F_p)=24\,N_{\mathrm{split}}(p)+12\bigl(2+n_c(p)\bigr)+6\,\varepsilon_\infty(p),
\qquad
\varepsilon_\infty(p)=\frac{1+\left(\frac{-1}{p}\right)}{2}.
\]
Here \(N_{\mathrm{split}}(p)\) counts the \(y\in\F_p\setminus\{0,1\}\) for which \(c(y)\neq 0\) and \(\Pi(\lambda,y)\) splits completely over~\(\F_p\), while
\[
n_c(p)=\#\{\beta\in\F_p:c(\beta)=0\text{ and }\Pi(\lambda,\beta)\text{ splits into linear factors over }\F_p\}.
\]
Equivalently, \(n_c(p)\) counts those rational roots \(\beta\) of~\(c\) for which the two simple roots of \(\Pi(\lambda,\beta)\) lie in~\(\F_p\), and \(\varepsilon_\infty(p)\) records whether \(-1\) is a square modulo~\(p\). The first fifteen values are:
\begin{center}
\input{sections/generated/prym_traces_table}
\end{center}
For every good prime \(p<120\), the companion computation verifies the Weil bound \(|a_1(Q,p)|\leq 6\sqrt{p}\), the divisibility
\[
3\mid a_1(X,p)-a_1(\Jac(X_A),p)-2\,a_1(E_{\mathrm{res}},p)-3\,a_1(E,p),
\]
and agreement with the second quotient-theoretic trace \(a_1(Y,p)-a_1(E_{\mathrm{res}},p)\), where \(Y=X/C_4\). These quotient computations check the subgroup bookkeeping and the resulting trace identities; the contribution above \(y=\infty\) is certified separately by the local expansion below.
\end{proposition}

\begin{proof}
Let \(y_0\in\PP^1(\F_p)\), and let \(D\) and \(I\) be the decomposition and inertia groups at a point of~\(X\) above~\(y_0\). The closed points above \(y_0\) are indexed by \(S_4/D\), each with residue degree \([D:I]\), while the geometric fibre has \([S_4:I]\) points. Thus the fibre has \([S_4:I]\) rational points when \(D=I\), and none otherwise.

At a finite branch point, \(I\) is generated by a transposition. Tameness gives \(D/I\subseteq N_{S_4}(I)/I\), a group of order~\(2\) generated by the complementary transposition. Hence \(D=I\) exactly when Frobenius fixes the two unramified sheets, or equivalently when \(\Pi(\lambda,y_0)\) splits into linear factors over~\(\F_p\). An unramified finite fibre contributes \(24\) rational points exactly when \(\Pi(\lambda,y_0)\) splits completely over~\(\F_p\), and none otherwise. The factorizations
\[
\Pi(\lambda,0)=\lambda(\lambda-1)^2(\lambda+1),\qquad
\Pi(\lambda,1)=(\lambda+1)^2(\lambda-1)(\lambda-2)
\]
show that the fibres at \(0\) and~\(1\) always contribute~\(12\). A rational root \(\beta\) of~\(c\) contributes~\(12\) precisely when the residual quadratic, equivalently the two simple roots of \(\Pi(\lambda,\beta)\), splits over~\(\F_p\). This gives the first two terms of the formula.

At infinity, set \(s=1/y\), \(t^4=s\), and \(\lambda=t^{-2}L\). Substitution into \(\Pi(\lambda,y)=0\), followed by multiplication by \(t^8\), gives
\[
L^4-t^2L^3-(2+t^4)L^2+t^6L+1+t^4=0.
\]
At \(t=0\) this is \((L^2-1)^2=0\). Substituting \(L=1+tZ\) and dividing by \(t^2\) gives \(4Z^2-1\) at \(t=0\), with simple roots \(Z=\pm\tfrac12\); substituting \(L=-1+tZ\) instead gives \(4Z^2+1\), with simple roots \(Z=\pm\tfrac{i}{2}\). Hensel's lemma therefore yields two branches over \(\QQ[[t]]\) and two over \(\QQ(i)[[t]]\) that do not descend to \(\QQ[[t]]\). Explicitly, the four branches begin
\[
L=1\pm\frac12t+\frac14t^2\pm\frac{7}{64}t^3+\cdots,
\qquad
L=-1\pm\frac{i}{2}t+\frac14t^2\mp\frac{7i}{64}t^3+\cdots.
\]
Two branches have coefficients in~\(\QQ\), while the other two require~\(i\). Hence the inertia group is \(I=C_4=\langle\tau\rangle\), the decomposition group over~\(\QQ\) is \(D=D_4\), and the residual quadratic extension is \(\QQ(i)/\QQ\). For odd~\(p\), the decomposition group after reduction is \(C_4\) when \(p\equiv1\pmod4\) and \(D_4\) when \(p\equiv3\pmod4\), so
\[
\#\pi^{-1}(\infty)(\F_p)=
\begin{cases}
6,&p\equiv1\pmod4,\\
0,&p\equiv3\pmod4.
\end{cases}
\]

By Theorem~\ref{thm:automorphic},
\[
3\,a_1(Q,p)=a_1(X,p)-a_1(\Jac(X_A),p)-2\,a_1(E_{\mathrm{res}},p)-3\,a_1(E,p).
\]
The traces \(a_1(E,p)\) and \(a_1(E_{\mathrm{res}},p)\) are obtained by point counting on their explicit Weierstrass models, while \(a_1(\Jac(X_A),p)\) is computed from \(w^2=-y(y-1)c(y)\). For every good prime below~\(120\), the script also applies the same decomposition-group count to \(E=X/S_3\), \(E_{\mathrm{res}}=X/D_4\), and \(X_A=X/A_4\), reproducing their model point counts, and to \(Y=X/C_4\), checking the subgroup bookkeeping and the resulting trace identities. These quotient computations test the finite-fibre criterion but not the contribution at infinity, which is certified by the local expansion and its dedicated symbolic assertion.
\end{proof}

\subsection{Conductor and reduction}

\begin{proposition}\label{prop:disc-quintic}
The discriminant of the hyperelliptic model \(f(y)=-y(y-1)c(y)\) of \(X_A\) is
\[
\Disc(f)=-2^{20}\cdot 3^{15}\cdot 31^2\cdot 37.
\]
The set of primes of bad reduction for \(X_A\) is contained in \(\{2,3,31,37\}\).
\end{proposition}

\begin{proof}
The roots of \(f(y)=-y(y-1)c(y)\) are \(0,1,\beta_1,\beta_2,\beta_3\), where \(\beta_i\) are the roots of \(c(y)\). For leading coefficient \(a_5=-256\),
\[
\Disc(f)=a_5^{8}\prod_{i<j}(r_i-r_j)^2.
\]
The product of pairwise differences factors as
\begin{align*}
\prod_{i<j}(r_i-r_j)^2
&=(0-1)^2\cdot\prod_i(0-\beta_i)^2\cdot\prod_i(1-\beta_i)^2\cdot\prod_{i<j}(\beta_i-\beta_j)^2.
\end{align*}
By Vieta's formulas for \(c(y)=256\prod_i(y-\beta_i)\):
\[
\prod_i\beta_i=-\frac{32}{256}=-\frac{1}{8},\qquad
\prod_i(1-\beta_i)=\frac{c(1)}{256}=\frac{864}{256}=\frac{27}{8}.
\]
Using \(\Disc(c)=256^4\prod_{i<j}(\beta_i-\beta_j)^2=-3^9\cdot 31^2\cdot 37\) (Proposition~\ref{prop:number-field}),
\[
\prod_{i<j}(r_i-r_j)^2=1\cdot\frac{1}{64}\cdot\frac{729}{64}\cdot\frac{-3^9\cdot 31^2\cdot 37}{256^4}
=\frac{-3^{15}\cdot 31^2\cdot 37}{2^{44}}.
\]
Therefore \(\Disc(f)=(-256)^8\cdot(-3^{15}\cdot 31^2\cdot 37)/2^{44}=2^{64}\cdot(-3^{15}\cdot 31^2\cdot 37)/2^{44}=-2^{20}\cdot 3^{15}\cdot 31^2\cdot 37\).
\end{proof}

\begin{theorem}\label{thm:reduction-types}
The reduction types of the genus-\(2\) curve \(X_A\) at the bad primes are as follows:
\begin{enumerate}
\item[(i)] At \(p=31\): the polynomial \(f(y)\bmod 31\) has exactly one double root (at \(y\equiv 12\)), giving semistable reduction with one ordinary double point. The conductor exponent is \(f_{31}=1\).
\item[(ii)] At \(p=37\): the polynomial \(f(y)\bmod 37\) has exactly one double root (at \(y\equiv 6\)), giving semistable reduction with one ordinary double point. The conductor exponent is \(f_{37}=1\).
\item[(iii)] At \(p=3\): the branch cubic reduces to \(c(y)\equiv (y-1)^3\pmod{3}\), so
\[
f(y)\equiv 2y(y-1)^4\pmod{3},
\]
and the special fiber has a non-split tacnode at \((1,0)\) of delta invariant~\(2\): the two local branches are conjugate over~\(\F_9\) and meet to second order. Moreover, the three non-rational branch points form a single depth-\(3/2\) cluster above \(y=1\), so \(X_A\) does not satisfy the semistability criterion of \cite[Definition~1.8 and Theorem~1.9]{DDMM2023}; in particular \(X_A/\QQ_3\) is not semistable.
\item[(iv)] The conductor of \(\Jac(X_A)\) has the form
\[
N_A=2^{f_2}\cdot 3^{f_3}\cdot 31\cdot 37,\qquad f_2\geq 1,
\]
and Theorem~\ref{thm:conductor-exact} identifies the wild part at \(3\) as \(\delta_3=1\).
\end{enumerate}
\end{theorem}

\begin{proof}
For~(i), one computes \(c(y)\equiv 8(y-6)(y-12)^2\pmod{31}\), so
\[
f(y)\equiv -8y(y-1)(y-6)(y-12)^2\pmod{31},
\]
with simple roots \(0,1,6\) and double root~\(12\). The double root produces a single ordinary node on the special fiber. For a tame prime with a single node, Liu's analysis \cite{Liu2002} gives \(f_p=1\).

For~(ii), modulo \(37\) the branch cubic factors as \(c(y)\equiv -3(y-6)^2(y-14)\pmod{37}\), giving
\[
f(y)\equiv 3y(y-1)(y-6)^2(y-14)\pmod{37},
\]
and the same argument gives \(f_{37}=1\).

For~(iii), the branch cubic satisfies \(c(y)\equiv (y-1)^3\pmod{3}\), so the special fiber of \(w^2=f(y)\) over~\(\F_3\) is \(w^2=2y(y-1)^4\). The singularity at \((1,0)\) is a non-split tacnode: substituting \(y=1+t\) gives \(w^2=2t^4(1+t)\), whose two local branches \(w=\pm\sqrt{2}\,t^2\sqrt{1+t}\) are conjugate over~\(\F_9\) (since \(\sqrt{2}\notin\F_3\)) and meet to second order, giving delta invariant~\(2\). Proposition~\ref{prop:cluster-pictures}(i) shows that each cubic branch point generates a totally ramified cubic extension of~\(\QQ_3\), while the full local splitting field has ramification index~\(6\), with wild inertia of order~\(3\). Since the semistability criterion of \cite[Definition~1.8]{DDMM2023} requires the root field \(K(\mathcal R)\) to have ramification degree at most~\(2\), Theorem~1.9 of loc.\ cit.\ implies that \(X_A/\QQ_3\) is not semistable.

For~(iv), modulo \(2\) the polynomial \(f\) reduces to \(y^2(y+1)^2\), which is a perfect square, so \(f_2\geq 1\). The wild part at \(3\) is identified in Theorem~\ref{thm:conductor-exact}.
\end{proof}

\begin{proposition}\label{prop:node-type}
The nodal singularities of the special fibers of \(X_A\) at the two semistable primes have opposite splitting types:
\begin{enumerate}
\item[(i)] At \(p=31\), the node is split: the tangent directions at the double root \(y\equiv 12\) are rational over~\(\F_{31}\). The normalization of the nodal special fiber is the elliptic curve \(v^2=-8y(y-1)(y-6)\) over~\(\F_{31}\).
\item[(ii)] At \(p=37\), the node is non-split: the tangent directions at the double root \(y\equiv 6\) are conjugate over~\(\F_{37^2}\). The normalization is the elliptic curve \(v^2=3y(y-1)(y-14)\) over~\(\F_{37}\).
\end{enumerate}
In both cases, the Jacobian \(\Jac(X_A)\) acquires toric rank~\(1\) at each prime, with the abelian part of the N\'eron model equal to the Jacobian of the normalization.
\end{proposition}

\begin{proof}
Write \(f(y)=(y-a)^2g(y)\) at the double root \(y\equiv a\pmod{p}\). The tangent cone of \(w^2=f(y)\) at the node is \(v^2=g(a)\) (after \(v=w/(y-a)\)), which splits over~\(\F_p\) if and only if \(g(a)\in(\F_p^\times)^2\).

At \(p=31\), \(a=12\): we have \(g(y)=-8y(y-1)(y-6)\) and
\[
g(12)\equiv -8\cdot 12\cdot 11\cdot 6\equiv 19\equiv 9^2\pmod{31},
\]
so the node is split.

At \(p=37\), \(a=6\): here \(g(y)=3y(y-1)(y-14)\) and
\[
g(6)\equiv 3\cdot 6\cdot 5\cdot(-8)\equiv 20\pmod{37}.
\]
The Legendre symbol \(\left(\frac{20}{37}\right)=\left(\frac{5}{37}\right)=-1\), so \(g(6)\) is a quadratic non-residue and the node is non-split.

In both cases \(v^2=g(y)\) is an elliptic curve over~\(\F_p\), and the single node gives toric rank~\(1\) in the N\'eron model \cite{Liu2002}.
\end{proof}

\begin{proposition}\label{prop:cluster-pictures}
The \(p\)-adic cluster pictures of the roots of \(f(y)=-y(y-1)c(y)\) at the two wild primes are as follows.
\begin{enumerate}
\item[(i)] At \(p=3\): the Newton polygon of the translated polynomial \(c(1+3t)/27=256t^3+393t^2+195t+32\) centered at the collapsed root \(y=1\) shows that the three roots \(\beta_i\) of \(c(y)\) satisfy \(v_3(\beta_i-1)=1\) and \(v_3(\beta_i-\beta_j)=3/2\) for all \(i\neq j\). The cluster picture is
\[
\{0\}\;\bigg|\;\underbrace{\{1\}\;\bigg|\;\underbrace{\{\beta_1,\beta_2,\beta_3\}}_{\text{depth }3/2}}_{\text{depth }1},
\]
where the roots \(\beta_i\) are conjugate over~\(\QQ_3\). Each root generates a totally ramified cubic extension of~\(\QQ_3\); the full local splitting field has ramification index~\(6\), with wild inertia of order~\(3\). Thus the wild inertia is nontrivial and \(\delta_3>0\).
\item[(ii)] At \(p=2\): the Newton polygon of \(c(y)=256y^3+411y^2+165y+32\) has vertices at \((0,v_2(32))\!=\!(0,5)\), \((1,v_2(165))\!=\!(1,0)\), \((2,v_2(411))\!=\!(2,0)\), \((3,v_2(256))\!=\!(3,8)\), giving slopes \(-5,0,8\). (With the convention that the Newton polygon slope between \((i,v_i)\) and \((i+1,v_{i+1})\) is \(v_{i+1}-v_i\), each root \(\beta\) satisfies \(v_2(\beta)=-\text{slope}\) of the corresponding segment.) Each segment has length~\(1\), so by Hensel's lemma all three roots lie in~\(\QQ_2\). Denote these roots \(\beta_1,\beta_2,\beta_3\). Since \(v_2(c(1))=5\) and \(v_2(c'(1))=0\), one has \(v_2(\beta_2-1)=5\). The cluster picture is
\[
\underbrace{\{0,\beta_1\}}_{\text{depth }5}\;\bigg|\;\underbrace{\{1,\beta_2\}}_{\text{depth }5}\;\bigg|\;\{\beta_3\},
\]
with two balanced clusters of depth~\(5\) and one isolated root of valuation \(-8\).
\end{enumerate}
\end{proposition}

\begin{proof}
For~(i), the substitution \(y=1+3t\) gives \(c(1+3t)=27(256t^3+393t^2+195t+32)=27h(t)\). A further substitution \(t=1+u\) yields
\[
h(1+u)=256u^3+1161u^2+1749u+876,
\]
with \(3\)-adic valuations \(v_3(876)=1\), \(v_3(1749)=1\), \(v_3(1161)=3\), \(v_3(256)=0\). The Newton polygon has a single segment of slope \(-1/3\) from \((0,1)\) to \((3,0)\), so all three roots satisfy \(v_3(t_i-1)=1/3\). In particular \(v_3(\beta_i-1)=1+v_3(t_i)=1\), and
\[
v_3(\beta_i-\beta_j)=1+v_3(t_i-t_j).
\]
The pairwise distances follow from \(\Disc(h)=-3^3\cdot 31^2\cdot 37\), giving
\[
\textstyle\sum_{i<j}v_3(t_i-t_j)=\tfrac12\bigl(v_3(\Disc(h))-4\cdot v_3(256)\bigr)=\tfrac32,
\]
and by symmetry each \(v_3(t_i-t_j)=1/2\), whence \(v_3(\beta_i-\beta_j)=3/2\). The Eisenstein computation in Proposition~\ref{prop:number-field} shows that each \(\beta_i\) generates a totally ramified cubic extension of~\(\QQ_3\). Since \(v_3(\Disc(c))=9\) is odd, the discriminant is not a square in~\(\QQ_3\); hence the local Galois group is \(S_3\), the splitting field has ramification index~\(6\), and its wild inertia has order~\(3\).

For~(ii), the \(2\)-adic Newton polygon of \(c(y)\) has vertices \((0,5)\), \((1,0)\), \((2,0)\), \((3,8)\), giving slopes \(-5\), \(0\), \(8\). The root \(\beta_1\) with \(v_2(\beta_1)=5\) clusters with~\(0\); the root \(\beta_2\) with \(v_2(\beta_2)=0\) satisfies \(v_2(\beta_2-1)=5\) by Newton's method applied to \(c(1+h)=0\), since \(v_2(c(1))=5\) and \(v_2(c'(1))=0\); the root \(\beta_3\) with \(v_2(\beta_3)=-8\) is isolated.
\end{proof}

\begin{proposition}\label{lem:local-three-prep}
Let \(F_w=\QQ_3(\beta_1)\), where \(\beta_1\) is any root of the branch cubic \(c(y)\). Then:
\begin{enumerate}
\item[(i)] \(F_w/\QQ_3\) is totally ramified of degree~\(3\);
\item[(ii)] the local discriminant exponent is \(v_3(\Delta_{F_w/\QQ_3})=3\);
\item[(iii)] the finite branch roots of \(f(y)=-y(y-1)c(y)\) form exactly three \(G_{\QQ_3}\)-orbits, namely
\[
\{0\},\qquad \{1\},\qquad \{\beta_1,\beta_2,\beta_3\}.
\]
\end{enumerate}
\end{proposition}

\begin{proof}
By the Eisenstein computation in Proposition~\ref{prop:number-field}, \([F_w:\QQ_3]=3\) and \(F_w/\QQ_3\) is totally ramified, proving~(i). Proposition~\ref{prop:cluster-pictures}(i) shows that the roots \(\beta_i\) are conjugate over~\(\QQ_3\) and satisfy \(v_3(\beta_i-1)=1\) and \(v_3(\beta_i-\beta_j)=3/2\). This cluster-theoretic description, together with the rationality of the roots \(0\) and~\(1\), proves the orbit statement~(iii).

For~(ii), Proposition~\ref{prop:number-field} gives \(\Disc(F/\QQ)=-3^3\cdot 37\). Since \(F_w/\QQ_3\) is totally ramified of degree~\(3\), the prime \(3\) has a unique extension to~\(F\) and residue degree~\(1\). Therefore the exponent of \(3\) in the global field discriminant equals the local discriminant exponent:
\[
v_3(\Disc(F/\QQ))=f_w\,v_3(\Delta_{F_w/\QQ_3})=1\cdot v_3(\Delta_{F_w/\QQ_3}),
\]
so \(v_3(\Delta_{F_w/\QQ_3})=3\).
\end{proof}

\begin{theorem}\label{thm:conductor-exact}
The wild part of the conductor of \(\Jac(X_A)\) at \(p=3\) is exactly
\[
\delta_3=1.
\]
\end{theorem}

\begin{proof}
Let
\[
R=\{0,1,\beta_1,\beta_2,\beta_3\}
\]
be the set of finite branch roots of the odd-degree model
\[
w^2=-256\prod_{r\in R}(y-r).
\]
The branch polynomial has odd degree \(5\), the residue characteristic is \(p=3\neq 2\), and the roots are distinct in \(\overline{\QQ}_3\) because \(\Disc(f)\neq 0\). Therefore the conductor formula of Dokchitser, Dokchitser, Maistret, and Morgan (Math. Ann. 385 (2023), Theorem~11.3) applies exactly in this setting: for an odd-degree hyperelliptic curve, the wild conductor is expressed as a sum over the \(G_{\QQ_3}\)-orbits on the finite branch set \(R\), with no extra contribution from the point at infinity. Since the scalar factor \(-1\) is a \(3\)-adic unit, it does not affect the extension fields \(\QQ_3(r)\) or the conductor contribution of any root. Thus
\[
\delta_3=\sum_{r\in R//G_{\QQ_3}}
\Bigl(v_3(\Delta_{\QQ_3(r)/\QQ_3})-[\QQ_3(r):\QQ_3]+f(\QQ_3(r)/\QQ_3)\Bigr),
\]
where \(R//G_{\QQ_3}\) denotes a set of \(G_{\QQ_3}\)-orbit representatives and \(f(\QQ_3(r)/\QQ_3)\) is the residue degree.

By Proposition~\ref{lem:local-three-prep}, the \(G_{\QQ_3}\)-orbits on \(R\) are \(\{0\}\), \(\{1\}\), and \(\{\beta_1,\beta_2,\beta_3\}\). The rational roots \(0\) and~\(1\) contribute
\[
v_3(\Delta_{\QQ_3/\QQ_3})-[\QQ_3:\QQ_3]+f(\QQ_3/\QQ_3)=0-1+1=0.
\]
Equivalently, the cluster picture contributes two rational singleton orbits and one cubic orbit centered \(3\)-adically at~\(1\), with pairwise distances \(v_3(\beta_i-\beta_j)=3/2\). All of the wild contribution is therefore concentrated in that single cubic orbit.
For the cubic orbit, Proposition~\ref{lem:local-three-prep} gives \([F_w:\QQ_3]=3\), \(f(F_w/\QQ_3)=1\), and \(v_3(\Delta_{F_w/\QQ_3})=3\).
Therefore the cubic orbit contributes
\[
v_3(\Delta_{F_w/\QQ_3})-[F_w:\QQ_3]+f(F_w/\QQ_3)=3-3+1=1,
\]
so the full sum is
\[
\delta_3=(0-1+1)+(0-1+1)+(3-3+1)=1.
\]
\end{proof}

\begin{remark}\label{rem:f2-obstruction}
The conductor exponent \(f_2\) remains undetermined. At \(p=2\), the special fiber of the naive hyperelliptic model is \(w^2\equiv y^2(y-1)^2\pmod{2}\), a non-reduced curve, so the minimal regular model over~\(\ZZ_2\) requires a nontrivial sequence of blowups whose combinatorics depend on the full \(2\)-adic cluster picture. Since all three roots of \(c(y)\) lie in~\(\QQ_2\) (Proposition~\ref{prop:cluster-pictures}(ii)), the stable reduction is in principle computable, but the two balanced clusters at depth~\(5\) make the computation substantially more involved than the wild-part calculation at~\(p=3\).
\end{remark}

\subsection{Torsion and Galois representations}

\begin{theorem}\label{thm:two-torsion}
The \(2\)-torsion field of \(\Jac(X_A)\) coincides with the splitting field of the branch cubic:
\[
\QQ(\Jac(X_A)[2])=L.
\]
The mod-\(2\) Galois representation
\[
\bar\rho_2\colon G_{\QQ}\longrightarrow \mathrm{Sp}(4,\F_2)\cong S_6
\]
has image isomorphic to \(S_3\), of order~\(6\) inside \(|\mathrm{Sp}(4,\F_2)|=720\).
\end{theorem}

\begin{proof}
For a genus-\(2\) curve \(C:w^2=f(y)\) with \(\deg f=5\), the nonzero elements of \(\Jac(C)[2]\) are represented by divisor classes of the form
\[
[(r_i,0)-(r_j,0)],
\]
where \(r_i,r_j\) are Weierstrass points; equivalently, the \(2\)-torsion field is the splitting field of~\(f\) \cite[Ch.~2]{CasselsFlynn1996}. In our case
\[
f(y)=-y(y-1)c(y)
\]
has Weierstrass points above \(0\), \(1\), and the three roots \(\beta_1,\beta_2,\beta_3\) of~\(c(y)\). Since \(0\) and~\(1\) are rational, adjoining all \(2\)-torsion points amounts to adjoining the cubic branch points \(\beta_i\). Hence
\[
\QQ(\Jac(X_A)[2])=\QQ(\beta_1,\beta_2,\beta_3)=L.
\]
Let \(W_0=\langle D_0,D_1\rangle\), where \(D_0=[(0,0)-\infty]\) and \(D_1=[(1,0)-\infty]\). Then \(W_0\) is fixed pointwise by \(G_{\QQ}\). In the quotient \(\Jac(X_A)[2]/W_0\), the classes \(\bar D_{\beta_1},\bar D_{\beta_2}\) form an \(\F_2\)-basis and satisfy \(\bar D_{\beta_3}=\bar D_{\beta_1}+\bar D_{\beta_2}\). Therefore a \(3\)-cycle \(\sigma=(123)\) and a transposition \(\tau=(12)\) act by
\[
\sigma\mapsto
\begin{pmatrix}
0&1\\
1&1
\end{pmatrix},
\qquad
\tau\mapsto
\begin{pmatrix}
0&1\\
1&0
\end{pmatrix}
\]
on \(\Jac(X_A)[2]/W_0\). These matrices generate \(\mathrm{GL}(2,\F_2)\cong S_3\), so the mod-\(2\) image is exactly \(S_3\). Since the Weil pairing is Galois invariant, this image embeds into \(\mathrm{Sp}(4,\F_2)\) as a subgroup of order~\(6\).
\end{proof}

\begin{corollary}\label{cor:two-torsion-classfield}
Theorem~\ref{thm:ray-class}(iv) identifies \(L\) as the unique cubic subextension of \(K_{\mathfrak p_3^2}/K\). The \(2\)-torsion of \(\Jac(X_A)\) is therefore controlled by the ray class group \(\Cl_{\mathfrak p_3^2}(K)\cong C_{24}\), and its absolute discriminant is \(\Disc(L/\QQ)=3^7\cdot 37^3\) (Theorem~\ref{thm:conductor}(iv)). The same extension~\(L/\QQ\) also determines the weight-one Artin representation of Corollary~\ref{cor:weight-one}.
\end{corollary}

\begin{proof}
Combine Theorem~\ref{thm:two-torsion} with Theorem~\ref{thm:ray-class}(iv).
\end{proof}

\subsection{The \texorpdfstring{$\F_2[S_3]$}{F2[S3]}-module structure of the \texorpdfstring{$2$}{2}-torsion}

Theorem~\ref{thm:two-torsion} identifies the \(2\)-torsion field as the splitting field~\(L\) of~\(c(y)\) and shows that the mod-\(2\) image is~\(S_3\). We now refine this to a complete description of the module structure and the resulting embedding.

\begin{theorem}\label{thm:mod2-structure}
As an \(\F_2[\Gal(L/\QQ)]\)-module, the \(2\)-torsion admits a decomposition
\[
\Jac(X_A)[2]\cong W_0\oplus W_1,
\]
where \(W_0=\langle D_0,D_1\rangle\cong \F_2^2\) is the trivial module spanned by the rational Weierstrass divisor classes and \(W_1\cong \mathrm{St}_2\) is the standard two-dimensional representation of~\(S_3\cong \mathrm{GL}(2,\F_2)\). Both \(W_0\) and \(W_1\) are nondegenerate (hyperbolic) subspaces for the Weil pairing, and the decomposition is a symplectic orthogonal sum:
\[
\bigl(\Jac(X_A)[2],e_2\bigr)\cong (W_0,e_2|_{W_0})\perp (W_1,e_2|_{W_1}).
\]
Consequently, the mod-\(2\) Galois image lies in the subgroup
\[
\bar\rho_2(G_{\QQ})\cong S_3\hookrightarrow \{1\}\times \mathrm{Sp}(W_1)\subset \mathrm{Sp}(W_0)\times \mathrm{Sp}(W_1)\subset \mathrm{Sp}(4,\F_2),
\]
acting trivially on the rational hyperbolic plane \(W_0\) and faithfully on the irrational hyperbolic plane~\(W_1\).
\end{theorem}

\begin{proof}
For the genus-\(2\) curve \(w^2=f(y)\) with \(\deg f=5\), the \(2\)-torsion \(\Jac(X_A)[2]\cong (\ZZ/2\ZZ)^4\) is generated by the divisor classes \(D_i=[(r_i,0)-\infty]\) for the five roots \(r_i\) of~\(f\), subject to \(\sum D_i=0\) \cite{CasselsFlynn1996}. The roots are \(0,1,\beta_1,\beta_2,\beta_3\). Taking \(\{D_0,D_1,D_{\beta_1},D_{\beta_2}\}\) as a basis, the Galois group \(S_3=\Gal(L/\QQ)\) fixes \(D_0\) and \(D_1\), and permutes \(\{D_{\beta_1},D_{\beta_2},D_{\beta_3}\}\) with \(D_{\beta_3}=D_0+D_1+D_{\beta_1}+D_{\beta_2}\).

For a genus-\(2\) Jacobian in a Weierstrass basis, the Weil pairing on \(2\)-torsion satisfies \(e_2(D_i,D_j)=-1\) for \(i\neq j\) \cite{CasselsFlynn1996}. Hence the Gram matrix in additive \(\F_2\)-notation with respect to \(\{D_0,D_1,D_{\beta_1},D_{\beta_2}\}\) is
\[
M=\begin{pmatrix}0&1&1&1\\1&0&1&1\\1&1&0&1\\1&1&1&0\end{pmatrix}.
\]
Consider the symplectic basis \(e_1=D_0,\;f_1=D_1,\;e_2=D_{\beta_1}+D_{\beta_2},\;f_2=D_0+D_1+D_{\beta_2}\). A direct computation gives
\[
\langle e_1,f_1\rangle=1,\quad \langle e_2,f_2\rangle=1,\quad \langle e_i,e_j\rangle=\langle f_i,f_j\rangle=\langle e_i,f_j\rangle=0\quad (i\neq j),
\]
exhibiting the standard symplectic form \(J_2\perp J_2\). The plane \(W_0=\langle e_1,f_1\rangle\) is fixed by~\(S_3\), and its symplectic complement \(W_1=\langle e_2,f_2\rangle\) carries the standard representation (the quotient action on \(\{D_{\beta_1},D_{\beta_2},D_{\beta_3}\}\) modulo~\(W_0\) is the standard \(2\)-dimensional representation of \(S_3\cong \mathrm{GL}(2,\F_2)\)). Both summands are hyperbolic, and the decomposition is \(S_3\)-equivariant.

The subgroup of \(\mathrm{Sp}(4,\F_2)\) preserving the orthogonal splitting \(W_0\perp W_1\) is \(\mathrm{Sp}(W_0)\times\mathrm{Sp}(W_1)\cong S_3\times S_3\). Since Theorem~\ref{thm:two-torsion} identifies the quotient action on \(W_1\) with the standard matrices above and \(S_3\) acts trivially on \(W_0\), the image is \(\{1\}\times S_3\).
\end{proof}

\begin{corollary}\label{cor:weight-one-mod2}
Let \(\rho_c\) be the standard two-dimensional Artin representation attached to \(L/\QQ\), and let \(\bar\rho_{2,W_1}\) be the action of \(G_{\QQ}\) on the nontrivial symplectic plane~\(W_1\). Both factor faithfully through \(\Gal(L/\QQ)\cong S_3\). For every prime \(p\) unramified in~\(L\), the cycle type of \(\operatorname{Frob}_p\) simultaneously determines the weight-one coefficient \(a_p(f_c)\) and the conjugacy class of \(\bar\rho_{2,W_1}(\operatorname{Frob}_p)\).
\end{corollary}

\begin{proof}
By Theorem~\ref{thm:hecke-traces}, \(a_p(f_c)\) is the standard Artin character evaluated at~\(\operatorname{Frob}_p\). Theorems~\ref{thm:two-torsion} and~\ref{thm:mod2-structure} identify \(\bar\rho_{2,W_1}\) with the faithful standard \(S_3\)-module. Thus the same Frobenius cycle type determines both assertions.
\end{proof}

\begin{remark}\label{rem:hyperbolic-splitting}
The plane \(W_0=\Jac(X_A)[2]^{G_{\QQ}}\) is the unique \(2\)-dimensional trivial \(S_3\)-submodule, since every trivial submodule lies in the fixed space. Its symplectic complement \(W_1=W_0^\perp\) is the unique \(S_3\)-stable complement: the set of such complements is a torsor under \(\Hom_{\F_2[S_3]}(W_1,W_0)\), which vanishes because \(W_1\) is the standard module and \(W_0\) is trivial. Hence the orthogonal decomposition \(W_0\perp W_1\) is unique. This rigidity contrasts with the situation for curves with fewer rational Weierstrass points, where the fixed space may be trivial and the Galois representation cannot be reduced in this manner.
\end{remark}

\subsection{Component groups at the semistable primes}

\begin{proposition}\label{prop:tamagawa}
Let \(\Phi_p\) denote the component group scheme \(\pi_0\) of the special fibre of the N\'eron model of \(\Jac(X_A)\) over~\(\F_p\), a finite \`etale \(\F_p\)-group scheme. Write \(\Phi_p(\overline{\F}_p)\) for its geometric component group with Frobenius action and \(\Phi_p(\F_p)\) for its group of rational components, so that \(c_p=\#\Phi_p(\F_p)\). Then:
\begin{enumerate}
\item[(i)] At \(p=31\), one has \(\Phi_{31}(\overline{\F}_{31})\cong\ZZ/2\ZZ\) with trivial Frobenius action, hence \(\Phi_{31}(\F_{31})\cong\ZZ/2\ZZ\) and \(c_{31}=2\).
\item[(ii)] At \(p=37\), the group scheme \(\Phi_{37}\) is trivial, hence \(\Phi_{37}(\overline{\F}_{37})=\Phi_{37}(\F_{37})=0\) and \(c_{37}=1\).
\end{enumerate}
\end{proposition}

\begin{proof}
At each semistable prime the special fibre has one ordinary double point and its normalization is an elliptic curve (Theorem~\ref{thm:reduction-types}). For a one-nodal fibre of thickness~\(n\), the dual graph of the minimal regular model is a cycle of length~\(n\), and Raynaud's description \cite[Chapter~10]{BoschLutkebohmertRaynaud1990} and \cite[Chapter~10]{Liu2002} gives
\[
\Phi_p(\overline{\F}_p)\cong\ZZ/n\ZZ.
\]
For a non-split node Frobenius acts on this cyclic group by~\(-1\), so in general \(c_p=\#(\ZZ/n\ZZ)[2]=\gcd(2,n)\). Thus non-splitness alone does not make the component group trivial.

The thickness is read from the colliding roots of \(f(y)=-y(y-1)c(y)\). By Proposition~\ref{prop:disc-quintic}, \(v_{31}(\Disc(f))=2\). The colliding roots lie in~\(\QQ_{31}\) at distance of valuation~\(1\), so the node has thickness \(n=2\). Blowing up once gives the two-component cycle in the minimal regular model. Proposition~\ref{prop:node-type} shows that the node is split, so Frobenius acts trivially on \(\Phi_{31}(\overline{\F}_{31})\cong\ZZ/2\ZZ\), proving~(i).

At~\(37\), one has \(v_{37}(\Disc(f))=1\); the colliding roots are conjugate over the ramified quadratic extension of~\(\QQ_{37}\), and the node has thickness \(n=1\). Hence \(\Phi_{37}=0\) as a group scheme, proving~(ii). The non-splitness from Proposition~\ref{prop:node-type} instead determines the toric part of the identity component: it is the non-split norm-one torus rather than~\(\mathbb G_m\).

For comparison, Proposition~\ref{prop:eres-conductor} gives \(\Delta_{\min}(E_{\mathrm{res}})=31^2\cdot37\). Its minimal \(c_4\) is a unit at both primes, so the reduction types are \(I_2\) at~\(31\) and \(I_1\) at~\(37\). Moreover, \(-c_6\equiv4\pmod{31}\) is a square, whereas \(-c_6\equiv8\pmod{37}\) is not. Thus \(E_{\mathrm{res}}\) has split \(I_2\) reduction with Tamagawa number~\(2\) at~\(31\), and \(I_1\) reduction with Tamagawa number~\(1\) at~\(37\), agreeing with \(\Jac(X_A)\) as stated in the summary.
\end{proof}

\subsection{Numerical verification}

\begin{proposition}\label{prop:verification}
The following table records the point counts \(\#X_A(\F_p)\) and \(\#E_{\mathrm{res}}(\F_p)\) at the first twelve primes of good reduction for~\(X_A\), together with the Frobenius traces and the number of roots of \(c(y)\) modulo~\(p\):
\[
\begin{array}{c|cccc|c}
p & \#X_A(\F_p) & a_1(X_A) & \#E_{\mathrm{res}}(\F_p) & a_1(E_{\mathrm{res}}) & \#\{y\in\F_p:c(y)=0\}\\ \hline
5 & 5 & 1 & 8 & -2 & 0 \\
7 & 5 & 3 & 13 & -5 & 0 \\
11 & 6 & 6 & 15 & -3 & 1 \\
13 & 14 & 0 & 16 & -2 & 1 \\
17 & 15 & 3 & 18 & 0 & 0 \\
19 & 18 & 2 & 26 & -6 & 1 \\
23 & 21 & 3 & 24 & 0 & 0 \\
29 & 31 & -1 & 34 & -4 & 0 \\
41 & 32 & 10 & 47 & -5 & 1 \\
43 & 48 & -4 & 38 & 6 & 1 \\
47 & 42 & 6 & 57 & -9 & 1 \\
53 & 62 & -8 & 59 & -5 & 1 \\
\end{array}
\]
The characteristic polynomial of Frobenius \(\chi_{13}(T)=T^4-2T^2+169\) is irreducible over~\(\QQ\), in agreement with Proposition~\ref{thm:simple}.
\end{proposition}

\begin{proof}
The point counts are computed by evaluating the Legendre symbol \((f(y)/p)\) for each \(y\in\F_p\), together with the point at infinity. The characteristic polynomial is determined by \(a_1=p+1-\#X_A(\F_p)\) and \(a_2=(a_1^2-p^2-1+\#X_A(\F_{p^2}))/2\). Irreducibility of \(\chi_{13}(T)=T^4-2T^2+169\) is shown in the proof of Proposition~\ref{thm:simple}.
\end{proof}

\begin{remark}\label{rem:verification-hecke}
The Hecke eigenvalues \(a_p(f_c)=\#\{y\in\F_p:c(y)=0\}-1\) of the weight-one form, computed from the last column of the table, match the values predicted by Theorem~\ref{thm:hecke-traces} at all primes listed. The values cycle through \(\{-1,0,2\}\), corresponding to the Frobenius cycle types \((123)\), \((12)\), and~\(1\) in~\(S_3\).
\end{remark}

\begin{remark}\label{rem:reproducibility-scripts}
For a structured inventory of the companion scripts and generated files used in the numerical checks of Sections~5 and~6, see the reproducibility note recorded after Section~7.
\end{remark}
